\documentclass{article}

\usepackage[margin=1.5cm]{geometry}
\usepackage[colorlinks]{hyperref}

\begin{document}

\section*{Essential Criteria}

\subsection*{Qualifications, experience and knowledge}

\paragraph{Possession of a PhD or equivalent qualification in a relevant subject
area}
%
I hold a PhD in Electrical Engineering from the
\href{https://www.usc.edu/}{University of Southern California}, completed
within its \href{https://minghsiehece.usc.edu/groups/sipi/}{Signal and Image
Processing Institute}.

\paragraph{An innovative personal research agenda in NeuroAI, computational
neuroscience, or a A/I closely related field.}
%
Please refer to my Research Vision described in the attached
rapelaVisionStmt.pdf.

\paragraph{Evidence of the ability to deliver and publish internationally excellent research.
Publications should demonstrate both methodological rigour (e.g. development or
application of quantitative/mathematical frameworks) and engagement with genuine
biological or neuroscientific questions, not purely technical contributions to AI/ML alone.}
%
I have followed a non-traditional yet highly rigorous career path uniquely
tailored to the modern demands of computational neuroscience. Following a PhD
and two postdoctoral fellowships at the interface of statistical machine
learning and neuroscience—resulting in key publications in leading journals
such as \textit{Neural Computation}, \textit{Network: Computation in Neural
Systems}, and the \textit{Journal of Vision}—I chose to embed myself deeply
within world-class experimental environments. I joined the Gatsby Computational
Neuroscience Unit and the Sainsbury Wellcome Centre (SWC) as a Research
Engineer Fellow to directly bridge the gap between advanced statistical
methodology and world-class neuroscience experimentation.

At the Gatsby Unit and SWC I have grown enormously. I have learned world class
machine learning at the Unit, and meeting one-on-one weekly with the Unit's
director, Maneesh Sahani, for more than seven years. I have acquired
professional software development skills. I have been part of a unique
computational/experimental neuroscience collaborative environment, where I
actively participated in the institute-wide team science Aeon project. I was
awarded £800.000 BBSRC award, and wrote a joint application between the Gatsby
Unit, the SWC and the Allen Institute for Neural Dynamics for a BBSRC/NSF-Bio
award. I lead a small team of research software engineers for more than three
years. I designed and taught the Statistical Neuroscience course for SWC PhD
students for two years. And I mentored two undergraduate students in
one-year-long fellowships sponsored by the Simmons Foundation.

Over the past seven years, my work has focused on enabling international,
institution-wide team science. Working closely with Unit Director Maneesh
Sahani, I have led a team of research software engineers, co-developed major
open-source ecosystems like \textit{Bonsai.ML}, and driven the statistical
infrastructure for the institute-wide \textit{Aeon} project. Crucially, I have
demonstrated the ability to secure major funding and deliver complex academic
milestones, including serving as a primary driver for an £800,000 BBSRC award
and a joint BBSRC/NSF-Bio application with the Allen Institute for Neural
Dynamics. Furthermore, I have committed to pedagogical excellence, designing
and teaching the foundational Statistical Neuroscience course for SWC PhD
students and mentoring Simmons Foundation fellows.

However, during this time I did not focus on research publications, which I
regret. Now I want to center my academic activities around publishing research
articles, teaching and mentoring students.
%
I am excellently prepared to do this. In the last seven years I have
accumulated papers ideas, and performed their background research, that will allow me
to publish several high-quality papers in the next few years. Below I summarise
a few of these paper ideas:

\begin{description}

    \item[traveling waves during rhythmic speech production] In Rapela 2016,
        2017, 2018, I reported, for the first time, these traveling waves in
        one subject. Now I have recordings from five more subjects and, if
        confirmed, I should be able to publish these finding in a high-impact
        journal.

    \item[svGPFA distribution paper] I will document the numerical,
        statistical and software engineering unique aspects in the Python
        distribution of svGPFA.

    \item[svGPFA scientific paper] Duncker and Sahani (2018) presented the
        method with only preliminary evidence of its potential to discover
        interesting structure in neural recordings. This paper will present
        robust svGPFA findings, maybe for detecting replay in the striatum.

    \item[svGPFA extensions] I will extend svGPFA to allow latent structure
        driven by external input and to operate in an online fashion.

    \item[Bonsai.ML paper] We will write a journal article presenting the
        Bonsai.ML package for neuroscience intelligent experimental control.

    \item[Aeon data processing paper] In collaboration with the undergraduate
        student that pioneered the statistical analysis of Aeon data, we will
        report methods and results from statistical analysis of Aeon data.

    \item[State-space models for behavioral inference] Current
        deep-neural-network models for pose estimation are not time-series
        models, and fail to use the temporal continuity of behavior. We have
        developed
        \href{https://joacorapela.github.io/ssm/auto_examples/index.html#kinematics-inference}{state-space
        models to infer behavioral kinematics}. We will describe these methods
        and highlight their benefits for pose estimation, of single and
        multiple animals.

\end{description}

\paragraph{Career Track Evolution: Deepening Team Science Through Publication.}

For the past seven years, my contributions to computational neuroscience have
been centered on institutional team science and methodological infrastructure.
Working in close collaboration with Gatsby Unit Director Maneesh Sahani, I have
focused on mastering world-class machine learning frameworks, leading
engineering teams, and co-developing foundational open-source ecosystems like
\textit{Bonsai.ML} and the \textit{Aeon} project. This path prioritized
large-scale tool creation, funding delivery, and pedagogical contributions over
a traditional personal publication track.

I am applying for this faculty position because I am executing a clear career
re-orientation regarding my role within team science. I want to transition away
from primary software development and engineering infrastructure. Instead, I
want to pivot my professional energy toward the scientific core of these
collaborations: formulating hypotheses, performing deep data analysis, and
leading the writing and publication of research papers. 

My goal is not to become an isolated investigator, but to expand the scope of
my collaborative work. I want to bring the elite machine learning methodology I
have developed at Gatsby into active, co-authored research programs with
experimental collaborators across NPP and CDB. This Lecturer/Associate
Professor role represents the ideal environment to make this transition,
allowing me to remain a dedicated team scientist while shifting my primary
output to peer-reviewed publications.

\paragraph{Career Re-orientation and Research Focus.}

For the past seven years, my academic contributions have been intentionally
centered on institutional team science and methodological infrastructure.
Working closely with Gatsby Unit Director Maneesh Sahani, I have focused my
efforts on mastering world-class machine learning frameworks, leading
engineering teams, and co-developing vital open-source ecosystems like
\textit{Bonsai.ML} and the \textit{Aeon} project. This path allowed me to
prioritize large-scale funding delivery—such as driving an £800,000 BBSRC
award—and pedagogical contributions over the traditional path of continuous
personal publication.

I am applying for this faculty position because I am making a distinct and
permanent career re-orientation. Having spent nearly a decade building,
scaling, and mastering the tools of computational neuroscience, my goal is to
transition away from primary software development. I want to shift my core
professional focus toward the independent formulation of research questions,
deep data analysis, and the consistent writing and publication of scientific
papers. 

While my recent track record reflects an engineer's focus on software systems
and tools, it has provided me with an elite methodological foundation. I am
seeking this Lecturer/Associate Professor role at UCL as the dedicated
environment in which to execute this transition, leveraging my deep expertise
in machine learning to establish a traditional, publication-focused research
program in NeuroAI.

\paragraph{Strategic Research Vision and Publication Velocity.} 

My career strategy has been deliberate and focused on long-term scientific
impact. Over the past seven years, I made a conscious choice to prioritize
institutional leadership, large-scale funding delivery, and the creation of
foundational software ecosystems like \textit{Bonsai.ML} and \textit{Aeon}.
While this path prioritized team science over an immediate volume of personal
publications, it was a necessary and highly rewarding investment. 

As I transition into a permanent faculty role, my core academic focus is
evolving. Having successfully established these world-class methodological and
software foundations, my primary objective is now the aggressive execution and
publication of my research pipeline. I am moving away from primary software
development to focus entirely on driving these mature, unique datasets into
peer-reviewed journals. The faculty position at UCL represents the exact moment
where my deep methodological preparation meets a fully realized, independent
publication program.

\paragraph{Evidence of the ability to deliver and publish internationally excellent research.}  

I have followed a non-traditional yet exceptionally rigorous career path
uniquely tailored to the modern demands of NeuroAI. Following a PhD and two
postdoctoral fellowships at the interface of statistical machine learning and
neuroscience, I chose to embed myself deeply within world-class experimental
and theoretical environments at UCL. For the past seven years, I have worked at
the absolute forefront of computational methodology through an intensive,
continuous intellectual partnership with Gatsby Unit Director Maneesh Sahani.
Meeting one-on-one weekly for over seven years, this collaboration has allowed
me to master and co-develop world-class machine learning frameworks explicitly
optimized for complex neural data.

Rather than treating machine learning as an isolated engineering tool, my
career has focused on enabling international, institution-wide team science
that bridges the gap between advanced mathematical theory and genuine
biological discovery. Leading a team of research software engineers across the
Gatsby Unit and the Sainsbury Wellcome Centre (SWC), I have co-developed major
open-source ecosystems like \textit{Bonsai.ML} and driven the statistical
infrastructure for the institute-wide \textit{Aeon} project. Crucially, I have
demonstrated the leadership required of a permanent faculty member by securing
funding and delivering complex academic milestones, serving as a primary driver
for an £800,000 BBSRC award and a joint BBSRC/NSF-Bio application with the
Allen Institute for Neural Dynamics. Furthermore, I have committed to
pedagogical excellence, designing and teaching the foundational Statistical
Neuroscience course for SWC PhD students and mentoring Simmons Foundation
fellows.

Having successfully established these world-class methodology, software, and
educational foundations, my research program is strategically positioned to
fully capitalize on this ecosystem. The unique insights, advanced
methodologies, and multi-subject datasets I have generated over the last seven
years have formed an immediate, high-impact publication pipeline. Below, I
outline the core manuscripts currently under preparation or extension that
directly marry world-class quantitative frameworks with fundamental
neuroscientific questions:

\begin{description}  
    \item[Traveling waves during rhythmic speech production:] Expanding on single-subject findings reported in my earlier work (Rapela 2016, 2017, 2018), I am leveraging a newly acquired dataset across six subjects to publish a robust characterization of speech-related cortical traveling waves.  
    \item[svGPFA distribution and scientific papers:] I am preparing a pair of manuscripts documenting the unique numerical and software engineering architecture of our Python distribution of svGPFA, alongside a scientific paper applying the method to detect complex latent structures, such as replay in the striatum.  
    \item[svGPFA extensions:] I am extending the core svGPFA framework to incorporate external latent drivers and real-time, online operation, moving the method beyond its current constraint to offline, preliminary evidence (Duncker and Sahani, 2018).  
    \item[Bonsai.ML and Aeon data processing papers:] I am leading the writing of journal articles that introduce the \textit{Bonsai.ML} package for intelligent experimental control and, in collaboration with an exceptional undergraduate mentee, reporting novel statistical insights from the \textit{Aeon} project.  
    \item[State-space models for behavioral inference:] To address the failure of deep pose-estimation networks to account for temporal continuity, we have developed state-space models to infer kinematics. I will describe these methods and demonstrate their advantages in capturing multi-animal behavioral continuity.  
\end{description}

\paragraph{Evidence of the ability to deliver and publish internationally excellent research.}  

I have followed a non-traditional yet highly rigorous career path uniquely tailored to the modern demands of computational neuroscience. Following a PhD and two postdoctoral fellowships at the interface of statistical machine learning and neuroscience—resulting in key publications in leading journals such as \textit{Neural Computation}, \textit{Network: Computation in Neural Systems}, and the \textit{Journal of Vision}—I chose to embed myself deeply within world-class experimental environments. I joined the Gatsby Computational Neuroscience Unit and the Sainsbury Wellcome Centre (SWC) as a Research Engineer Fellow to directly bridge the gap between advanced mathematical frameworks and genuine biological discovery.

Over the past seven years, my work has focused on enabling international, institution-wide team science. Working closely with Unit Director Maneesh Sahani, I have led a team of research software engineers, co-developed major open-source ecosystems like \textit{Bonsai.ML}, and driven the statistical infrastructure for the institute-wide \textit{Aeon} project. Crucially, I have demonstrated the ability to secure major funding and deliver complex academic milestones, including serving as a primary driver for an £800,000 BBSRC award and a joint BBSRC/NSF-Bio application with the Allen Institute for Neural Dynamics. Furthermore, I have committed to pedagogical excellence, designing and teaching the foundational Statistical Neuroscience course for SWC PhD students and mentoring Simmons Foundation fellows.

Having successfully established these world-class software, data, and educational foundations, my career is strategically transitioning to fully capitalize on this ecosystem. The unique insights, advanced methodologies, and multi-subject datasets I have generated over the last seven years have formed an immediate, high-impact publication pipeline. Below, I outline the core manuscripts currently under preparation or extension that directly marry rigorous quantitative frameworks with fundamental neuroscientific questions:

\begin{description}  
    \item[Traveling waves during rhythmic speech production:] Expanding on single-subject findings reported in my earlier work (Rapela 2016, 2017, 2018), I am leveraging a newly acquired dataset across six subjects to publish a robust, high-impact characterization of speech-related cortical traveling waves.  
    \item[svGPFA distribution and scientific papers:] I am preparing a pair of manuscripts documenting the unique numerical and software engineering architecture of our Python distribution of svGPFA, alongside a scientific paper applying the method to detect complex latent structures, such as replay in the striatum.  
    \item[svGPFA extensions:] I am extending the core svGPFA framework to incorporate external latent drivers and real-time, online operation, moving the method beyond its current constraint to offline, preliminary evidence (Duncker and Sahani, 2018).  
    \item[Bonsai.ML and Aeon data processing papers:] I am leading the writing of journal articles that introduce the \textit{Bonsai.ML} package for intelligent experimental control and, in collaboration with an exceptional undergraduate mentee, reporting novel statistical insights from the \textit{Aeon} project.  
    \item[State-space models for behavioral inference:] To address the failure of deep pose-estimation networks to account for temporal continuity, we have developed state-space models to infer kinematics. I will describe these methods and demonstrate their advantages in capturing multi-animal behavioral continuity.  
\end{description}

\paragraph{Evidence of potential to secure competitive research funding.}

\paragraph{Evidence of e ective supervision and mentoring of undergraduate,
postgraduate and/or doctoral research students including teaching in a
computational, data science, or quantitative context (e.g. programming,
machine learning, statistics, mathematical modelling) — ideally in neuroscience
or bioscience settings; Python experience strongly preferred.}

\paragraph{Experience of research-based teaching at undergraduate and/or postgraduate levels,
including willingness to gain Fellowship of the Higher Education Academy if equivalency A/I
not already obtained.}

\paragraph{Evidence of productive cross-disciplinary collaborations,
specifically spanning computational or AI methods and neuroscience or biological
systems, for example joint A/I publications, grants, or collaborative projects
with researchers from complementary disciplines.}

\paragraph{Experience of engagement beyond academia, including contributions to
policy, practice, public engagement, or knowledge exchange.}

\subsection*{Skills and abilities}

\paragraph{Ability to develop and implement team-based solutions to complex challenges,
delivering positive outcomes.}

\paragraph{Advanced computational skills relevant to neuroscience research, including proficiency
in at least one scientific programming language (e.g. Python, MATLAB, R, Julia) and
familiarity with machine learning frameworks or data analysis tools used in neuroscience
or bioscience.}

\subsection*{UCL Ways of Working}

\paragraph{Ability to build collaborative research and industry partnerships.}

\paragraph{Ability to develop and implement team-based solutions to complex challenges,
delivering positive outcomes.}

\paragraph{Active commitment to equality, diversity, and inclusion, for example
through inclusive teaching practices, mentoring underrepresented students,
participating in EDI initiatives, I or fostering an inclusive lab/group
culture.}

\section*{Desirable Criteria}

\paragraph{Postgraduate Certificate in Teaching \& Learning in Higher Education, or equivalent
qualification.}

\paragraph{Experience of contributions to peer review bodies/committees, professional
organisations, learned societies, government committees or Research Councils etc.}

\paragraph{Experience of an active role in widening participation initiatives
that span research communities, e.g. sharing ideas and experience through
organising seminars and workshops.}

\end{document}
